\section{Discussion}
\label{sec:discussion}

\paragraph{What the random split measures}
The results are consistent with a single explanation. Fraud occurs in episodes, in which a compromised card is used several times within minutes or days and the transactions of an episode resemble each other. A random split distributes each episode over the training and test data, so the test measures whether a model can recognize the remainder of an episode that it has already seen labeled. This task is much easier than detecting a new episode, and a random split additionally provides the model with every other transaction of the episode, including later ones. The chronological split poses the question that a deployed system faces. The difference between the two is not primarily gradual concept drift, which would grow with the distance from the training period, because nearly all of it is present in the first week.

The size of the gap depends on how much a model can memorize. Logistic regression loses the least in absolute terms on IEEE-CIS and shows no measurable loss from the oversampling leak, whereas the tree ensembles and the MLP lose the most. A random split therefore overstates performance and also overstates the benefit of model capacity.

\paragraph{What the two datasets can support}
IEEE-CIS is large enough for these effects to be measured precisely, but ULB is not. Its \num{ulb.fraud} fraud cases over two days leave a chronological test set of \num{ulb.testfraud} cases and random-split scores that vary by several points from seed to seed, so an improvement of a point or two of \prauc{} on a single split of ULB lies within the noise, whereas the oversampling leak is far larger than the noise.

\paragraph{A checklist for evaluating fraud models}
\begin{enumerate}[leftmargin=1.8em,itemsep=2pt]
\item \textbf{Split by time.} Train on earlier transactions and test on later ones, and state the cut points. Tune on a validation period that lies between them.
\item \textbf{Respect label delay.} Leave a gap between training and test periods that matches the time labels take to arrive.
\item \textbf{Fit all preprocessing on training rows only.} This includes imputation, scaling, encoding, feature selection and resampling. Resample after splitting, and never resample the test set.
\item \textbf{Score at the real class ratio.} Report \prauc{} and recall at a stated alert budget on unmodified test data. Accuracy and F1 on a rebalanced test set are not informative.
\item \textbf{Check that a gap is leakage.} A later period can be harder. Compare protocols on the same transactions, as in Section~\ref{sec:controls}, before attributing a difference to the split.
\item \textbf{Repeat over periods and report variance.} Use rolling windows where the data allow it. On small datasets, state how many fraud cases the test set contains and how much the score varies across splits.
\item \textbf{Report value-weighted metrics.} Recall by value can differ from recall by count.
\item \textbf{Release code and stored predictions,} so that others can recompute metrics under a different protocol.
\end{enumerate}

\paragraph{Limitations}
Both datasets are public and anonymized, and we deliberately used raw features. A production system would add card-history aggregates, which change both the level of performance and possibly the size of the gap, so our numbers are not estimates of production performance. The card key we use on IEEE-CIS is a heuristic, and we do not know how or when that dataset's labels were assigned. For reasons of cost, the IEEE-CIS oversampling experiment used a 200{,}000-transaction subsample with three seeds, and the rolling windows and cross-validation used one seed. The chronological protocols score a fixed model on six weeks of data and do not model retraining or investigator feedback, which the realistic-evaluation literature treats in depth~\cite{dalpozzolo2018credit}.

\ifanon\else
\begin{acks}
In preparing this paper, the authors used Claude Code (Anthropic) as a coding and writing assistant. We used it to implement the experiment pipeline, run the experiments, produce drafts of the text, and check references. We designed the study, specified the experiments and the analyses, reviewed and edited all code, results and text, and take full responsibility for the content.
\end{acks}
\fi

\section{Conclusion}

On the larger of the two public payment-fraud datasets, the \prauc{} of five standard classifiers is \num{ieee.min.rel.pr}--\num{ieee.max.rel.pr}\% lower on future transactions than a random split suggests, and oversampling before the split inflates scores further. The results suggest that overlap between fraud episodes in the training and test data is an important contributor to this inflation in both cases. On the smaller dataset, the split matters less than it first appears and the oversampling leak matters more, and the dataset is too small to support the fine-grained model comparisons that are often drawn from a single split of it. Time-ordered splits, resampling confined to the training data, alert-budget metrics and a comparison on the same transactions are inexpensive to adopt, and they make a benchmark score a closer estimate of the performance a fraud team would observe.
